\documentclass[10pt,a4paper]{amsart}
\usepackage[margin=19mm]{geometry}
\usepackage{amsmath,amssymb,mathtools}
\usepackage[scr=rsfs,cal=euler]{mathalfa}
\usepackage{xfrac}
\usepackage[unicode,hidelinks,bookmarks=false]{hyperref}
\newcommand{\ie}{i.e.\ }
\newcommand{\cf}{cf.\ }
\newcommand{\resp}{respectively\ }
\newcommand{\Subsection}{\subsection}
\providecommand{\nobreakdash}{\nobreak\hbox{-}}
\setlength{\emergencystretch}{2em}
\multlinegap0pt
\usepackage{amssymb}
\newtheorem{lemma}{Lemma}
\title{On the Existence of Balanced Generalized de Bruijn Sequences}
\author{Matthew Baker, Bhumika Mittal, Haran Mouli,, Eric Tang}
\date{Source: arXiv:2201.11863v1 (CC BY 4.0)}
\begin{document}
\maketitle

\noindent\emph{Source and licence.} On the Existence of Balanced Generalized de Bruijn Sequences. Version arXiv:2201.11863v1 (CC BY 4.0), distributed by arXiv under the Creative Commons Attribution 4.0 International licence, http://creativecommons.org/licenses/by/4.0/, as recorded in the arXiv licence metadata for this exact version and shown on the abstract page https://arxiv.org/abs/2201.11863v1. Full licence notice: \texttt{licenses/arxiv-2201.11863-CC-BY-4.0.txt}.\par\smallskip

\noindent\textit{Editorial scope.} Counted unit: the article's lemma nin (source0.tex lines 199--201). The article's complete original proof of it is delivered with it (source0.tex lines 203--213, one \texttt{proof} environment of 11 lines). No mathematics is written, repaired or re-derived: the statement and the proof are the article's own lines, in the article's own notation, and no retained line is substituted or re-flowed.  Retained uncounted as the source-local context the proof needs: lines 165--167; lines 189--191.  Nothing was omitted from any retained span, so the package carries no omission record.  No retained line contains a citation, so the wrapper carries no bibliography and no bibliography entry.  No retained line contains a figure environment, an image-inclusion command or drawing code, so no image is delivered and the package declares no asset dependency.  Editorial formatting only, in addition: an AMS \texttt{amsart} preamble that restates the article's own declarations and macros used by the retained lines, namely the environments \texttt{lemma} on separate counters, together with the standard packages those lines need.  The retained environments print in this document as Lemma 1 in order of appearance, whereas the article prints the same items with the numbering of its own full document.  The complete native source of the article as distributed in the arXiv e-print for this exact version is bundled unchanged as \texttt{sources/121314/source0.tex}.
\begin{lemma}\label{redblue}
Each vertex $v \in V(G_l)$ has one red and one blue out-edge, and the two in-edges of $v$ have the same color.
\end{lemma}

\begin{lemma}\label{cycle}
If the graph $G_l$ has a balanced circuit of length $n$, then the graph $G_{l+1}$ has a balanced cycle of length $n$.
\end{lemma}

\begin{lemma}\label{nin}
Let $l \geqslant 2$, and let $n$ be an even positive integer at most $2^l$. Then the graph $G_l$ contains a balanced circuit of length $n$.
\end{lemma}

\begin{proof}
We proceed by induction on $l$. First, for $l=2$, we can either have $n=2$ or $n=4$. We may take the circuits $0 \to 1$ and $0 \to 1 \to 1 \to 0$, respectively. 

Now, assume that the claim holds true for $l$. We prove it holds true for $l+1$. For even $n \leqslant 2^l$, we may use the induction hypothesis to find a balanced circuit of length $n$ in $G_l$. From Lemma \ref{cycle}, there exists a balanced cycle (also a circuit) of length $n$ in $G_{l+1}$. 

We now show the claim for $2^l < n < 2^{l+1}$ where $n$ is even. Since $2 \leqslant 2^{l+1}-n < 2^l$, there exists a balanced $n$-cycle in $G_{l+1}$, say $H$. Then the graph $G_{l+1}-H$ contains $n$ edges. Assume that $G_{l+1}-H = H_1 \cup H_2 \cup \cdots \cup H_t$, where each $H_i$ for $1 \leqslant i \leqslant t$ is a component of $G_{l+1}-H$. Since both $G_{l+1}$ and $H$ are balanced and Eulerian, $G_{l+1}-H$ must be balanced and each $H_i$ must be Eulerian.

If $G_{l+1}-H$ is connected, we are done, since it is a circuit in $G_{l+1}$ of length $n$. Otherwise, since $G_{l+1}$ is connected, there exists an edge $e$ in $H$ which connects vertices in two different components of $G_{l+1}-H$, say $H_1$ and $H_2$. Let the edge $e$ go from $v_1 \in V(H_1)$ to $v_2 \in V(H_2)$. Without loss of generality, assume $e$ is red. Now, let $e_1$ be an edge from $v_1$ to $u_1$ in $H_1$, and let $e_2$ be an edge from $u_2$ to $v_2$ in $H_2$. Then the last $l-1$ bits of $v_1$ and $u_2$ are the first $l-1$ bits of $v_2$ and $u_1$. Hence, there must be an edge $e'$ from $u_2$ to $u_1$ in $H$. We deduce from Lemma \ref{redblue} that since $e$ is red, $e_1$ must be blue, $e_2$ must be red, and $e'$ must be blue.

Consider the new subgraph obtained by replacing $e_1$ and $e_2$ with $e$ and $e'$. One checks easily that the degrees of all vertices, as well as the number of red and blue edges, are preserved. Moreover, the number of connected components has been reduced. Thus, we end up with a balanced subgraph with $n$ edges and fewer components, all of which are Eulerian. By repeating this process, we will eventually end up with a balanced subgraph of $n$ edges which is a circuit. Finally, for $n=2^{l+1}$, all of the edges of the graph form a circuit since $G_{l+1}$ is Eulerian. The result follows.
\end{proof}

\end{document}
